\section{制度背景}\label{sec:institution}

\subsection{价格工具：一个三十年未动的税目}

现行消费税对“鞭炮、焰火”按 \textbf{15\% 从价}征收，纳税环节为生产、委托加工与进口。以下四项事实构成本文制度背景的骨架，其证据档如下。

\begin{table}[H]
\centering
\caption{价格工具的四项制度事实与证据档}
\label{tab:excise-facts}
\small
\begin{tabular}{L{10.9cm}L{4.1cm}}
\toprule
\textbf{事实} & \textbf{证据档} \\
\midrule
税目“鞭炮、焰火”，从价 15\%，生产／委托加工／进口环节缴纳 & 一手：国家税务总局 12366 口径与现行税目税率表 \\
该税目为 1993-12-31《消费税暂行条例》原始 11 个税目之一，1994-01-01 起征；\textbf{不是} 2006 年新增 & 一手 \\
2006-04-01 财税〔2006〕33 号新增木制一次性筷子（5\%）与实木地板（5\%），\textbf{未改动该税目} & 一手 \\
1994 年该税目税率同为 15\% & 两份独立转载互证；\tk{1993 年附件原件尚未从政府站点调取} \\
\bottomrule
\end{tabular}
\end{table}

这里必须纠正一个常见的误记：该税目\textbf{不是} 2006 年那一轮“绿色化”调整的产物。它比那一轮早十二年，而且那一轮并没有碰它。1993 年设置该税目时，应税消费品被分类，\textbf{第一类即“过度消费危害人类健康、社会秩序和生态环境的特殊消费品——如烟、酒、鞭炮、焰火”}，理由写作“寓禁于征”。也就是说，\textbf{环保理由从一开始就在，而且排在第一位}——这一点与“这个税当初只是为了收钱”的直觉相反。\tk{1993 年官方释义原件；现有表述来自教科书与转载}

这个税目还有两个结构特征，比税率本身更值得注意。

\textbf{一是出口免征。} 消费税对出口应税消费品免征，税基是国内销售。中国是全球最大的烟花爆竹生产国与出口国，出口量以十万吨计（2024 年约 40.8 万吨）。\textbf{占产量相当大一块的部分，完全不在这一税种的射程之内。}

\textbf{二是从价而非从量。} 15\% 按销售额征收，而排放正比于火药质量。同样药量的廉价爆竹与高价礼花，税负可以相差数倍，污染却一样。\textbf{每单位污染对应的税负随产品结构任意变动。} 这比任何关于税率的讨论都更能说明：这个税在设计上就不是对着污染量去的。

\begin{quote}
这两点合起来支持一个判断：\textbf{该税目在设立时带有环境理由，但其征收结构从未按环境目标设计过。} 它是一条挂在一般消费税征管上的税目，不是一个环境工具。
\end{quote}

\subsection{数量工具：从上级要求到地方自主加码}\label{sec:ban-history}

禁放不是一开始就由地方自主决定的。2017 年的国家大气污染防治方案（环大气〔2017〕110 号）第 25 项要求“各个城市要制定烟花爆竹禁放限放严控方案，明确春节期间限放区域和允许燃放时间，有条件的城市建成区内全时段禁止销售、燃放烟花爆竹”。那一波禁放带有明确的上级要求色彩。

2023 年之后情况变了。本文对 34 个城市的大气污染防治方案做了全文取证，并与两份国家级文件对照：

\begin{itemize}
  \item 《空气质量持续改善行动计划》（国发〔2023〕24 号，39 项）\textbf{全文不含“烟花爆竹”}；
  \item 《京津冀及周边地区、汾渭平原 2023—2024 年秋冬季大气污染综合治理攻坚方案》（环大气〔2023〕73 号）\textbf{全文不含“烟花爆竹”}。
\end{itemize}

在这两份国家级文件均为零命中的背景下，34 个城市的表现是：

\begin{table}[H]
\centering
\caption{34 城大气污染防治方案中“烟花爆竹”的出现情况}
\label{tab:34cities}
\small
\begin{tabular}{L{7.8cm}L{7.2cm}}
\toprule
\textbf{口径} & \textbf{结果} \\
\midrule
方案中提到烟花爆竹管控（含只写半句话） & 23 城（占取得原文的 30 城中的 77\%） \\
写成独立编号任务 & 11 城（占 30 城中的 37\%） \\
官方全文零命中（确证没有） & 7 城：潍坊、太原、沈阳、鞍山、丹东、昆明、杭州 \\
文件未找到（既不能算有，也不能算没有） & 4 城：滨州、广州、东莞、绵阳 \\
\bottomrule
\end{tabular}
\end{table}

\noindent 注：若把 4 个未取得原文的城市计入分母，则“提到”的比例为 23/34 $\approx$ 68\%。两种口径都在附录~\ref{app:documents} 中列明；正文采用“取得原文的城市”作分母，与“未找到不等于没有”的纪律一致。

\noindent \textbf{一处必须交代的限定。} 23 个“提到”里有 5 个来自 2023 年之前的文件（周口 2020、佛山 2017、石家庄 2021、武汉 2022、宝鸡 2022），它们可能反映的是 2017 年国家方案的要求，而不是 2023 年之后的地方自主选择。只看 \textbf{2023 年及以后印发}的 25 份可判定文件，仍有 \textbf{18 份提到（72\%）、8 份独立列项（32\%）}——与全样本比例接近，因此结论不是靠旧文件撑起来的。

\noindent \textbf{样本性质。} 34 城是\textbf{便利样本}（山东河南 11、华南华东 11、华北 12），不是随机抽样。上述比例只描述这 34 城，\textbf{不能推广为“全国城市方案”}。

\textbf{这就是本节的关键事实：2023 年之后，城市方案里写不写这条，是地方自主加码，不是上级摊派。} 同为“空气质量持续改善”类文件，太原、沈阳、鞍山、丹东四市完全没有，而西安把它升格为独立专项。

省级框架仍然塑造市级表述，这一点有两处一手对照：河北省级文件（冀政发〔2024〕4 号，26 项）中“烟花爆竹”只出现 1 次、是“秸秆禁烧”条的末句，唐山与石家庄的表述与之高度同构；湖北省级文件（鄂政发〔2024〕6 号，10 个专项行动）中只出现 1 次、位于“面源精细化”之内，武汉与宜昌的一句话结构与之同构。\textbf{上级不强制，但上级的措辞被下级沿用。}

\subsection{一条完整的制度链路}

上述论证容易被视为“文件里写了什么”而无实质。东营提供了一个反例：一条从环境考核文件到政府规章、再到实际收紧的完整链路。

《东营市 2023—2024 年秋冬季大气污染综合治理攻坚行动方案》（东营市生态环境委员会办公室 2023 年 11 月 16 日印发）把烟花爆竹列为 7 个专项行动之第（六）项，其第 18 项原文：

\begin{quote}
\textbf{（六）烟花爆竹燃放管控专项行动}

\textbf{18.} 2023 年 12 月底前，制定完成《东营市烟花爆竹燃放管理办法》，做好烟花爆竹禁限放立法衔接工作。严格落实烟花爆竹燃放相关管理规定，在禁止燃放区域、时段内实施禁燃禁放……（牵头单位：市公安局、市应急局、市司法局、市生态环境局）
\end{quote}

那部《办法》就是 2024 年 1 月 2 日施行的东营市人民政府令第 208 号。其第十条规定“本市下列区域内\textbf{全年禁止}燃放烟花爆竹”，全文检索“除夕／正月初一／正月十五／限制燃放”均为 0 次——即取消了此前《暂行办法》（2021-09-22）给予的 7 天燃放豁免。

\textbf{大气攻坚方案 $\rightarrow$ 市级政府规章 $\rightarrow$ 2024 年实际收紧}，一条链完整闭合。这是本项目找到的唯一一份“环境考核文件直接下令制定禁令规章”的一手证据。它同时说明了禁令的制度位置：写在以生态环境部门为主的方案里、与量化工业减排指标同文下达；东营把市公安局列在牵头单位首位，但牵头单位各地不一（见下）。

\begin{quote}
\textbf{必须说明的一处边界}：牵头单位各地不一。河南一侧由应急管理居首，合肥与宜昌由市生态环境局牵头，佛山由市民族宗教局牵头，南京 2024 版、宁波、青岛、南阳、周口、唐山、大连等根本没有写牵头单位。\textbf{不得写成“由公安局牵头”。}
\end{quote}
